\documentclass[UTF8,12pt,a4paper,fontset=fandol]{ctexart}

% 支持从项目根目录、深度学习目录或当前文件目录加载共享样式。
\makeatletter
\def\input@path{{../}{../../}}
\makeatother
\usepackage{researchnotes}

\title{变换器原理与面试问答}
\author{}
\date{2026年10月3日}

\begin{document}
\maketitle
\tableofcontents

\section{研究对象、学习路线与统一记号}
\label{sec:tr-intro}

Transformer 通常保留英文名称，也常译为\keyterm{变换器}。它是一类以注意力机制组织序列信息交换的神经网络架构，而不是一种特定训练目标、分词算法或文本生成策略。理解它需要同时掌握三个层次：一个词元如何读取其他词元的信息，多个计算模块如何组成深层网络，以及网络如何通过训练形成条件概率并在推理时生成结果。本文先从可手算的注意力出发，建立矩阵公式与张量形状，再讨论位置、掩码、残差、归一化和前馈网络，随后说明训练、缓存与计算代价，最后以分类问答检验理解。原始架构的历史来源是文献\cite{vaswani2017}；后续变体分别注明来源，不把某一种模型配置视为所有 Transformer 的定义。

贯穿例子是续写“我今天想喝”。模型接收的实际对象是分词器产生的\keyterm{词元（token）}编号，词元可能是一个汉字、子词、标点或若干字节，不能预设上述文本一定恰好对应五个词元。假设教学分词结果为五个词元，则模型先把五个编号映射为向量，经过若干层信息交换，最后用第五个位置的输出给“茶”“水”等候选词元分配概率。一次前向传播可以计算多个已知位置的预测，但若后续词元尚未生成，就不能把它们当作已知输入。这一区别贯穿训练并行性与推理串行性的讨论。

本文面向掌握矩阵乘法、链式法则和基本概率的读者，问题覆盖概念解释、公式推导、维度核算、代码实现、训练排错和推理系统。所谓“常见面试题”指这些基础能力所对应的代表性问题，并非基于招聘频率调查得出的排名。回答时应先说明结论，再给公式或例子，最后交代条件；例如“注意力复杂度是二次的”必须指出相对于序列长度、哪一部分计算，以及是否包含投影和前馈网络。

设批量大小为 $B$，查询序列长度为 $n$，键值序列长度为 $m$，模型宽度为 $d$，层数为 $L$，查询头数为 $h$，单头查询与键维度为 $d_k$，单头值维度为 $d_v$，前馈中间宽度为 $d_f$，词表大小为 $U$。常见标准多头配置满足 $d_k=d_v=d/h$，但注意力公式本身不要求 $d_k=d_v$。单个序列的表示矩阵为 $X\in\mathbb R^{n\times d}$，各行对应位置；批量输入为 $X_{\rm batch}\in\mathbb R^{B\times n\times d}$。讨论单层公式时省略批量维，向量默认按行排列；涉及旋转矩阵时会显式改用列向量。矩阵 $V$ 专指值矩阵，避免同时用它表示词表大小。所有对数默认是自然对数，参数集合统一记为 $\theta$。

\section{从内容检索到缩放点积注意力}
\label{sec:tr-attention}

\subsection{查询、键和值分别解决什么问题}

\keyterm{注意力（attention）}可以理解为可微分的内容检索。查询（query）表达当前位置需要什么信息，键（key）提供每个候选位置用于匹配的特征，值（value）提供实际被汇总的内容。检索条件和返回内容没有必要相同：查图书时用主题匹配，取回的却可能是整段正文。这只是功能类比，真实的查询、键和值都是训练得到的连续向量，并没有预先指定某个坐标一定表示语法、实体或情感。

在自注意力中，令
\begin{equation}
Q=XW_Q,\qquad K=XW_K,\qquad V=XW_V,
\label{eq:tr-qkv}
\end{equation}
其中 $W_Q,W_K\in\mathbb R^{d\times d_k}$，$W_V\in\mathbb R^{d\times d_v}$，因此 $Q,K\in\mathbb R^{n\times d_k}$，$V\in\mathbb R^{n\times d_v}$。这里省略偏置以突出主要结构；具体模型可以保留或删除偏置。三个矩阵来自同一个 $X$，并不意味着三个投影矩阵相同。不同投影允许模型分别学习“怎样发问”“怎样被找到”和“提供什么内容”。直接规定 $Q=K=V=X$ 仍能定义一种注意力运算，但会限制可学习的匹配关系。

\keyterm{自注意力（self-attention）}的查询与键值来自同一序列；\keyterm{交叉注意力（cross-attention）}的查询来自一个序列，键和值来自另一个序列。例如翻译时，目标前缀产生查询，源句编码结果产生键值，此时通常 $n\ne m$。对一般输入 $Q\in\mathbb R^{n\times d_k}$、$K\in\mathbb R^{m\times d_k}$ 和 $V\in\mathbb R^{m\times d_v}$，缩放点积注意力为
\begin{equation}
S=\frac{QK^{\mathsf T}}{\sqrt{d_k}}+M,\qquad
A_{ij}=\frac{\exp(S_{ij})}{\sum_{r=1}^{m}\exp(S_{ir})},\qquad
O=AV\in\mathbb R^{n\times d_v}.
\label{eq:tr-attention}
\end{equation}
$M\in(\mathbb R\cup\{-\infty\})^{n\times m}$ 是加性掩码或偏置；禁止读取的位置设为 $-\infty$，无额外偏置的允许位置设为零。每个查询至少需要一个允许访问的键。Softmax 沿键的维度逐行计算，$A_{ij}$ 表示查询位置 $i$ 分配给键位置 $j$ 的权重，不是沿查询维归一化。

\begin{example}[一次完整的注意力计算]
取 $d_k=2$，只有一个查询 $q=(\sqrt2,0)$，三个键依次为 $k_1=(1,0)$、$k_2=(0,1)$、$k_3=(1,1)$，对应值为 $v_1=(1,0)$、$v_2=(0,2)$、$v_3=(3,1)$，且所有键可见。缩放后的分数为 $(1,0,1)$，故
\[
a=\frac{(e,1,e)}{2e+1}\approx(0.4223,0.1554,0.4223),
\qquad
o=\left(\frac{4e}{2e+1},\frac{2+e}{2e+1}\right)
\approx(1.6893,0.7330).
\]
如果第三个键被屏蔽，则必须重新归一化，得到 $a=(e/(e+1),1/(e+1),0)$，输出约为 $(0.7311,0.5379)$。把原来的第三个权重直接改成零却不重新归一化，会得到另一个运算。这个例子同时说明：输出由值决定，查询和键只负责权重；某个值可以有负坐标，注意力并不要求输出各坐标非负。
\end{example}

\subsection{为什么除以维度的平方根}

缩放因子的直觉是控制匹配分数的尺度，而不是把向量变成单位向量。设随机查询分量 $Q_r$ 和随机键分量 $K_r$ 均为零均值、单位方差，所有这些分量彼此独立，则
\[
\mathbb E[Q_rK_r]=0,\qquad
\operatorname{Var}(Q_rK_r)=\mathbb E[Q_r^2]\mathbb E[K_r^2]=1.
\]
不同坐标的乘积也独立，因此
\begin{equation}
\operatorname{Var}\left(\sum_{r=1}^{d_k}Q_rK_r\right)=d_k,
\qquad
\operatorname{Var}\left(\frac{\sum_{r=1}^{d_k}Q_rK_r}{\sqrt{d_k}}\right)=1.
\label{eq:tr-scale}
\end{equation}
这是初始化尺度分析的理想模型，不是训练后每个头必然满足的统计恒等式。若分量存在相关性，方差还含协方差项；若方差不是一，也需要相应调整分析。除以 $d_k$ 会使上述模型中的分数方差变成 $1/d_k$，维度越大越趋于平坦；除以 $\sqrt{d_k}$ 则保持量级稳定。

设一行权重为 $a_j=\exp(s_j)/\sum_r\exp(s_r)$，则
\begin{equation}
\frac{\partial a_j}{\partial s_r}=a_j(\delta_{jr}-a_r),
\label{eq:tr-softmax-grad}
\end{equation}
其中 $\delta_{jr}$ 为克罗内克符号。当分数差距很大、权重近似独热时，这个雅可比矩阵的许多元素接近零，注意力分配对分数的局部变化可能不敏感。缩放有助于避免初始化时过早进入这种状态，但不能据此断言所有损失梯度都会消失，例如输出 Softmax 与交叉熵合并求导有不同的梯度表达。缩放也没有保证注意力必然均匀；模型仍可通过学习形成尖锐的分配。

数值实现应使用 $a_j=\exp(s_j-c)/\sum_r\exp(s_r-c)$，其中 $c=\max_r s_r$。减去同一个有限常数不改变结果，却可避免指数溢出。如果整行都是 $-\infty$，则 $c=-\infty$，减法会产生未定义数值；这不是简单更换精度就能解决的问题，而是需要规定无有效键的处理方式。把禁止项写成零同样错误，因为 $\exp(0)=1$；把它写成较大的有限负数则需要确认在使用的数据类型及分数范围内确实足够小。

\subsection{权重能说明什么，不能说明什么}

当没有对权重施加随机失活时，$A_{ij}\geq0$ 且每行和为一，所以每个单头输出是相应值向量的凸组合。但经过输出投影、残差相加或后续前馈网络后，完整层的输出不再受这个凸包限制。训练中的注意力权重随机失活通常会对保留项重缩放；单次随机实现的行和不必等于一，因此凸组合结论也不能不加条件地套用。

即使 $Q=K$ 使原始点积分数矩阵对称，逐行 Softmax 后的 $A$ 通常仍不对称，因为不同的行有不同的归一化分母。一般情况下 $W_Q\ne W_K$，原始分数就未必对称。注意力权重反映特定层、特定头、给定输入下的信息混合系数，并不自动等于词元对最终答案的因果贡献；值的大小、其他头、残差分支和后续层都影响输出。因此展示热力图可以辅助检查，但不能只凭某个词权重最大就认定模型“真正依据”它作答。

\section{多头机制、前馈网络与深层计算}
\label{sec:tr-block}

\subsection{多头不是把一个结果重复计算多次}

\keyterm{多头注意力（multi-head attention, MHA）}使用多组投影，让同一个位置拥有多种匹配与汇总方式。对第 $a$ 个头，$a=1,\ldots,h$，记
\[
H_a=\operatorname{Attention}(XW_Q^{(a)},XW_K^{(a)},XW_V^{(a)}),
\]
则
\begin{equation}
\operatorname{MHA}(X)=\operatorname{Concat}(H_1,\ldots,H_h)W_O,
\qquad W_O\in\mathbb R^{hd_v\times d}.
\label{eq:tr-mha}
\end{equation}
每个头分别对自己的分数做 Softmax，之后才拼接。因此多个较窄的头通常不等价于一个同样总宽度的大头：一个大头只有一套位置权重，多头允许不同子空间采用不同权重。各头是否学到互补关系是训练结果，不能预先保证某个头专门负责主语、某个头专门负责指代。

当 $hd_k=hd_v=d$ 时，可将所有查询投影合并为一次 $XW_Q$，其结果仍为 $n\times d$，再按头拆分。批量实现常从 $(B,n,d)$ 变形为 $(B,n,h,d_k)$，交换轴后成为 $(B,h,n,d_k)$，以便最后两个维度做矩阵乘法。转置与重排不能混淆：仅改变形状可能把不同位置的数据混入同一个头。输出经过 $(B,h,n,d_v)$ 到 $(B,n,hd_v)$ 的逆过程，再乘 $W_O$，才能与原输入进行残差相加。

\subsection{逐位置前馈与门控变体}

\keyterm{逐位置前馈网络（position-wise feed-forward network, FFN）}对每个位置使用同一组参数，一种基本形式为
\begin{equation}
\operatorname{FFN}(X)=\phi(XW_1+\mathbf1b_1)W_2+\mathbf1b_2,
\quad W_1\in\mathbb R^{d\times d_f},\quad W_2\in\mathbb R^{d_f\times d}.
\label{eq:tr-ffn}
\end{equation}
这里 $\phi$ 是逐元素非线性，$b_1,b_2$ 为行偏置，$\mathbf1$ 是长度为 $n$ 的全一列向量。注意力直接混合不同位置的信息，FFN 在当前位置混合通道并实施非线性变换。虽然 FFN 本身不跨位置读取输入，但它处理的表示已可包含上文信息，因此不能说 FFN “只理解孤立词”。如果两层线性变换之间没有非线性，它们可合并为一层仿射变换，表达能力会受限。

常用非线性包括修正线性单元（rectified linear unit, ReLU）、高斯误差线性单元（Gaussian error linear unit, GELU）以及门控线性单元（gated linear unit, GLU）的变体。ReLU 为 $\max(0,x)$；GELU 的精确定义为 $x\Phi(x)$，其中 $\Phi$ 是标准正态分布函数，实际实现也可能使用近似式。令 $\operatorname{SiLU}(x)=x\sigma(x)$，$\sigma(x)=1/(1+e^{-x})$，一种不含偏置的 SwiGLU 写作
\begin{equation}
\operatorname{SwiGLU}(X)=
\bigl[\operatorname{SiLU}(XW_g)\odot(XW_u)\bigr]W_d,
\label{eq:tr-swiglu}
\end{equation}
其中 $W_g,W_u\in\mathbb R^{d\times d_f}$，$W_d\in\mathbb R^{d_f\times d}$，$\odot$ 表示逐元素乘积。门控给不同通道提供随输入变化的调制，不应把 SiLU 门值误认为严格位于零与一之间。相关设计见文献\cite{shazeer2020}。

忽略偏置时，普通 FFN 有 $2dd_f$ 个参数，SwiGLU 有 $3dd_f$ 个参数。因此比较激活函数时应说明是否同时改变中间宽度。若普通 FFN 取 $d_f=4d$，其参数数为 $8d^2$；门控版本若要大致匹配预算，可取 $d_f\approx8d/3$，实际常再按实现要求取整。模型参数更多和非线性设计更好是两种不同解释，公平实验需要控制这种差异。

\subsection{残差、层归一化与均方根归一化}

\keyterm{残差连接（residual connection）}写为 $Y=X+F(X)$，使子层学习对当前表示的修正，并提供直接的梯度通路。其雅可比为 $I+J_F(X)$，确实包含恒等项，但很多层相乘后仍可能出现梯度爆炸、消失或抵消，不能把残差当作绝对稳定性的证明。残差还要求两个分支形状相同，这解释了为什么注意力输出投影和 FFN 的第二个投影都回到宽度 $d$。

对一个位置的向量 $x\in\mathbb R^d$，\keyterm{层归一化（layer normalization, LayerNorm）}定义为
\begin{equation}
\mu=\frac1d\sum_{r=1}^d x_r,\qquad
v=\frac1d\sum_{r=1}^d(x_r-\mu)^2,\qquad
\operatorname{LN}(x)_r=\gamma_r\frac{x_r-\mu}{\sqrt{v+\varepsilon}}+\beta_r,
\label{eq:tr-ln}
\end{equation}
其中 $\gamma,\beta\in\mathbb R^d$ 可学习，$\varepsilon>0$ 防止除零。统计量在单个词元的特征维上计算，常规用法不混合不同样本，也不混合未来位置。归一化前的标准化部分在忽略 $\varepsilon$ 时具有零均值、单位方差；加入可学习缩放与平移后的最终输出一般不再具有这两个数值性质。

\keyterm{均方根归一化（root mean square normalization, RMSNorm）}的一种常见形式为
\begin{equation}
\operatorname{RMSNorm}(x)_r=
\gamma_r\frac{x_r}{\sqrt{d^{-1}\sum_{s=1}^d x_s^2+\varepsilon}}.
\label{eq:tr-rms}
\end{equation}
它用均方根控制尺度，不减去均值，常见实现也不带平移参数。它保留了不同于 LayerNorm 的均值信息，不能描述成“与 LayerNorm 完全等价，只是计算更快”。其提出和实验背景见文献\cite{zhang2019}；具体速度仍取决于内核融合、张量形状和硬件。

对于同一个子层 $F$，\keyterm{后归一化（Post-LN）}为 $Y=\operatorname{LN}(X+F(X))$，\keyterm{前归一化（Pre-LN）}为 $Y=X+F(\operatorname{LN}(X))$。原始 Transformer 使用前一种组织方式；许多后续架构采用后一种并在堆叠末尾再做归一化。Pre-LN 的残差主路不必每层穿过归一化雅可比，这有助于解释某些训练稳定性差异，但不能推出它在任意深度、学习率和初始化下都不需要预热。文献\cite{xiong2020}给出特定理论假设和实验配置下的分析，应保留这些适用边界。

忽略随机失活，一层采用 Pre-LN 的仅解码器块可写为 $\widetilde X=X+\operatorname{MHA}(\operatorname{LN}(X))$，$Y=\widetilde X+\operatorname{FFN}(\operatorname{LN}(\widetilde X))$，其中 MHA 内部使用因果掩码。使用 RMSNorm 或门控 FFN 时替换相应组件即可。\keyterm{随机失活（dropout）}在训练时引入随机掩码，推理通常关闭；不同配置可能在注意力权重、投影输出或残差分支使用它，不能仅凭“Transformer”名称假定所有位置都有相同失活率。

\section{位置表示与注意力可见范围}
\label{sec:tr-position}

\subsection{为何需要位置机制}

若没有位置表示和位置相关掩码，注意力只根据内容比较位置。设 $P$ 是置换矩阵，同时重新排列输入行，则 $Q'=PQ$、$K'=PK$、$V'=PV$，分数变为 $Q'K'^{\mathsf T}=P(QK^{\mathsf T})P^{\mathsf T}$。逐行 Softmax 与这个置换相容，所以
\begin{equation}
\operatorname{Attention}(PX)=P\operatorname{Attention}(X).
\label{eq:tr-permutation}
\end{equation}
这里将固定投影包含在 Attention 的记号中。这叫\keyterm{置换等变性（permutation equivariance）}：输入换序，输出随之换序；它不等于输出完全不变的置换不变性。逐位置 FFN、归一化和残差也保留这种性质，故仅堆叠这些模块无法让同一内容集合的不同顺序拥有独立的顺序语义。若最后再采用对称池化，才可能进一步得到整体输出的置换不变性。

上述证明明确排除了位置相关掩码。固定的因果三角掩码并不会在任意置换下保持不变，它本身提供顺序结构；因此“没有显式位置编码的所有 Transformer 都完全不知道顺序”是过强结论。不过只有因果可见范围仍可能不足以有效表达精细距离和长程规律，是否移除位置机制需要在具体任务和训练方案中检验。

\subsection{绝对位置、相对位置与旋转位置}

\keyterm{绝对位置编码（absolute positional encoding）}把位置编号映射为向量 $p_i$，与词元嵌入相加。可学习位置表为每个训练范围内的位置保留参数，超出表长需要扩展或重新设计。正弦位置编码在偶数宽度 $d$ 下使用
\begin{equation}
p_{i,2r}=\sin(i\omega_r),\qquad
p_{i,2r+1}=\cos(i\omega_r),\qquad
\omega_r=10000^{-2r/d},\quad 0\le r<d/2.
\label{eq:tr-sinusoidal}
\end{equation}
不同频率覆盖不同变化尺度，固定位置差对应每对正余弦坐标上的固定线性变换。公式可计算任意整数位置，不表示训练后模型一定能可靠处理任意长度；可定义性、数值可表示性和任务泛化能力是三个不同问题。

\keyterm{相对位置偏置（relative position bias）}直接在分数中加入依赖 $i-j$ 的项，例如 $S_{ij}=q_i k_j^{\mathsf T}/\sqrt{d_k}+b_{i-j}$。可以对距离分桶，以减少参数或表达远近层次。ALiBi，即注意力线性偏置（attention with linear biases），在因果场景的一种表达为对允许的 $j\le i$ 加入 $-c_a(i-j)$，其中第 $a$ 个头的斜率 $c_a>0$。它给更远位置更大的距离惩罚，而不是把远处内容直接删除；其长序列实验见文献\cite{press2021}，不能由偏置形式推导所有任务都能无限外推。

\keyterm{旋转位置编码（rotary position embedding, RoPE）}在查询和键的二维坐标对上施加与位置有关的旋转。为说明代数性质，此处将单个查询和键视为列向量。设
\[
R(\alpha)=\begin{pmatrix}\cos\alpha&-\sin\alpha\\\sin\alpha&\cos\alpha\end{pmatrix},
\qquad \widetilde q_i=R(i\omega)q_i,\quad \widetilde k_j=R(j\omega)k_j.
\]
因为 $R(\alpha)^{\mathsf T}R(\beta)=R(\beta-\alpha)$，有
\begin{equation}
\widetilde q_i^{\mathsf T}\widetilde k_j
=q_i^{\mathsf T}R((j-i)\omega)k_j.
\label{eq:tr-rope}
\end{equation}
高维情形由多个不同频率的二维旋转块组成。这解释了绝对位置旋转怎样把显式的相对位置差引入点积。公式不意味着整个分数只取决于距离：内容向量 $q_i,k_j$ 仍在式中，深层内容也可能已包含位置信息。旋转保持向量范数，且常规用法主要旋转查询和键；是否处理值向量取决于具体变体。方法来源见文献\cite{su2021}。

RoPE 的频率范围、位置索引、旋转配对方式和缓存中的位置偏移必须一致。把两个相邻坐标配对与把向量前后两半配对可以构造等价的不同实现约定，但权重排列也必须一致，不能只替换一段代码而默认兼容。调整频率、位置缩放或训练长度都可能帮助扩展上下文，却同时改变模型见到的位置分布，需要重新评估长距离检索、局部精度与不同长度上的稳定性。

\subsection{因果掩码、填充掩码与损失掩码}

\keyterm{因果掩码（causal mask）}要求位置 $i$ 只能读取 $j\le i$：允许项为零，未来项为 $-\infty$。对角线通常允许，因为输入位置 $i$ 的词元用于预测下一词元，而不是预测自身；防止答案泄漏还依赖第\ref{sec:tr-training}节的标签移位。\keyterm{填充掩码（padding mask）}禁止读取批处理占位符，和因果掩码解决不同问题。两者应共同决定每一行有效键的集合，屏蔽填充键并不自动把填充查询的输出清零。

\keyterm{损失掩码（loss mask）}决定哪些预测位置参与训练目标，不改变前向可见范围。指令微调中可以只对回答词元计算损失，但回答仍需读取提示词；若同时屏蔽提示词的键，模型就失去了必要条件。把多个样本拼成一条序列时，是否允许后一个样本读取前一个样本是数据建模约定：若要保持样本独立，应设置相应的块对角因果可见范围，必要时同步处理位置编号；仅加入结束符或忽略某些标签并不会在计算图中强制隔离样本。

增量解码时查询长度和键长度可能不同。若缓存已有 $p$ 个位置，新输入有 $q$ 个位置，则查询对应绝对位置 $p,p+1,\ldots,p+q-1$，键对应 $0,1,\ldots,p+q-1$，第 $r$ 个新查询允许键索引不超过 $p+r$。不能不加偏移地使用一个普通 $q\times(p+q)$ 左上三角矩阵，否则单词元查询可能只看见第一个键。不同框架对布尔掩码的真假语义也可能相反，必须查看对应接口并用极小例子检验。

\section{架构类型、训练目标与生成过程}
\label{sec:tr-training}

\subsection{编码器、解码器与编码器--解码器}

\keyterm{仅编码器（encoder-only）}模型通常在有效输入范围内采用双向注意力，使每个位置同时结合左右文，常用于分类、序列标注与表示学习。\keyterm{仅解码器（decoder-only）}模型通常在目标序列上采用因果注意力，并预测下一个词元。\keyterm{编码器--解码器（encoder--decoder）}模型先表示完整源序列，再由目标前缀通过交叉注意力读取源表示。这里的“解码器”指网络模块，\keyterm{解码策略（decoding strategy）}则指从输出分布选择词元的算法，二者不可混同。

原始翻译架构的编码器层包含自注意力和 FFN，解码器层另有读取编码结果的交叉注意力。仅解码器模型一般保留因果自注意力和 FFN，不要求存在一个未画出的编码器。BERT 是双向编码器预训练的代表，原论文包含掩码语言建模与下一句预测目标\cite{devlin2019}；后者并非所有编码器模型的必要组成。T5 采用统一文本到文本的编码器--解码器框架，并研究跨度破坏等训练目标\cite{raffel2020}。架构限制信息如何流动，训练目标规定学习什么，两者有关联，却不是一一对应关系。

\subsection{自回归似然、标签移位与有效词元平均}

设真实词元序列为 $x_1,\ldots,x_T$，并设起始符为 $x_0$。\keyterm{自回归语言建模（autoregressive language modeling）}采用概率链式分解
\begin{equation}
p_\theta(x_1,\ldots,x_T\mid x_0)
=\prod_{t=1}^T p_\theta(x_t\mid x_0,\ldots,x_{t-1}).
\label{eq:tr-autoreg}
\end{equation}
链式法则对任意联合分布成立，模型结构则约束条件概率如何被参数化和计算。将用于预测 $x_t$ 的最后一层表示记为 $h_t\in\mathbb R^d$，它经词表投影 $W_{\rm vocab}\in\mathbb R^{d\times U}$ 得到逻辑值（logits）$z_t=h_tW_{\rm vocab}+b$，随后 Softmax 输出词表分布。输入嵌入矩阵 $E\in\mathbb R^{U\times d}$ 有时与输出投影绑定，即 $W_{\rm vocab}=E^{\mathsf T}$，这减少参数，但不是所有配置的强制要求。

教学序列 $[\mathrm{BOS},\text{我},\text{爱},\text{茶},\mathrm{EOS}]$ 可构造输入 $[\mathrm{BOS},\text{我},\text{爱},\text{茶}]$、标签 $[\text{我},\text{爱},\text{茶},\mathrm{EOS}]$。也可以把完整序列交给内部执行移位的模型接口，但这时调用者不能再重复移位。训练利用真实前缀，称为\keyterm{教师强制（teacher forcing）}。输入各位置已知，所以在一层内部可以同时计算所有位置的注意力；因果掩码确保位置 $t$ 的预测不读取它应预测的答案。模型深度方向仍逐层依赖，参数更新也仍依赖反向传播。

令 $w_t\in\{0,1\}$ 表示位置是否参与损失，且 $N_{\rm valid}=\sum_t w_t>0$，则一次批处理的平均负对数似然为
\begin{equation}
\mathcal L(\theta)=-\frac{1}{N_{\rm valid}}
\sum_t w_t\log p_\theta(x_t\mid x_{<t}),
\qquad \operatorname{PPL}=\exp(\mathcal L).
\label{eq:tr-loss}
\end{equation}
批量公式把样本与位置共同求和即可。\keyterm{困惑度（perplexity, PPL）}是有效目标词元负对数似然的指数，不是生成答案的正确率。不同分词器会改变一个文本所包含的词元数量及预测事件，词元级困惑度通常不能直接横向比较；至少需要统一文本、分词、上下文与计分协议。跨长度样本求平均也要区分每词元等权和每样本等权，它们对应不同目标。

对于单位置的真实标签分布 $y$ 和预测概率 $p$，交叉熵对逻辑值的梯度为 $\partial\ell/\partial z=p-y$，可由 Softmax 雅可比与对数链式法则推得。实际应把未归一化逻辑值交给数值稳定的交叉熵实现，而不是先转概率再交给期望逻辑值的接口。标签平滑可把独热目标换成与背景分布的混合分布，缓和过度尖锐的训练目标，但会改变最优概率拟合对象，不能未经实验就断言必定提高准确率或校准性。

\subsection{训练配置与生成策略}

反向传播通过注意力、前馈网络、归一化与嵌入更新参数。AdamW、学习率预热、衰减、梯度裁剪和混合精度是训练方法，不是 Transformer 定义的一部分。预热让初始更新步幅逐渐增加，可以配合初始化与归一化缓解早期不稳定；有效性取决于具体规模与数据。\keyterm{混合精度（mixed precision）}在不同操作中采用不同数值精度，归约、归一化或优化器状态可能保留较高精度。FP16 动态范围较窄，损失缩放常用于避免梯度下溢；BF16 的指数范围较宽，是否需要损失缩放不能直接照搬 FP16 经验。

\keyterm{梯度累积（gradient accumulation）}把若干微批次的梯度合并后更新。在不同微批次有效目标数不相等时，简单平均各自的平均损失会让较短微批次权重过大；应按有效词元数还原全局求和与归一化。\keyterm{激活检查点（activation checkpointing）}减少保存的中间激活，在反向时重算部分前向，用额外计算换显存。它与 KV 缓存服务的阶段不同：前者主要降低训练反向所需激活存储，后者主要减少自回归推理时重复计算。

推理时先对提示前缀前向计算，再从最后一个有效位置的分布中选择下一词元并追加。\keyterm{贪心解码（greedy decoding）}每步取最大概率词元，只保证当前条件下的局部选择，不保证整个序列概率最大。\keyterm{束搜索（beam search）}保留多个候选前缀，近似搜索高分序列，但有限束宽仍无全局最优保证，且长度归一化会改变评分目标。开放式文本任务中，最高概率文本也不必最有信息或最符合用户期望。

\keyterm{温度（temperature）} $\tau>0$ 把采样分布改为 $p_i(\tau)\propto\exp(z_i/\tau)$。低温通常使分布更集中，高温使其更平坦；$\tau\to0^+$ 时质量集中到最大逻辑值集合，若有并列最大则需要说明平局规则。\keyterm{Top-$k$ 采样}只保留概率最大的 $k$ 个候选，\keyterm{核采样（nucleus sampling, Top-$p$）}保留按概率降序累积达到阈值的一组候选，两者都需对保留集合重新归一化。它们改变生成分布，不能保证避免错误或重复。

教师强制在真实前缀上训练，生成却会遇到模型自己产生的前缀，这种分布差异常称为\keyterm{暴露偏差（exposure bias）}。它说明局部预测误差可能影响后续输入，但不是所有生成错误的唯一来源；知识不足、目标不匹配、提示歧义和评测污染同样需要区分。长回答中的连贯性也不能仅靠调低温度保证，必须观察错误在不同前缀与长度上的传播。

\section{参数、计算量、显存与推理优化}
\label{sec:tr-efficiency}

\subsection{先明确核算边界}

在标准 MHA 中取 $hd_k=hd_v=d$，忽略偏置，查询、键、值和输出四个投影共有 $4d^2$ 个参数。普通 FFN 取 $d_f=4d$ 时有 $8d^2$ 个参数，因此一层仅解码器块的主要矩阵参数约为 $12d^2$；层归一化参数为 $O(d)$，通常是较小项。总参数还包括约 $Ud$ 的输入嵌入、是否独立的词表输出投影、位置参数及其他模块，编码器--解码器还要计入交叉注意力，不能把 $12Ld^2$ 当作任意 Transformer 的精确总参数。

以一个长为 $n$ 的序列为例，四个投影需要约 $4nd^2$ 次乘加，$QK^{\mathsf T}$ 和 $AV$ 合计约 $2n^2d$ 次乘加，普通 FFN 需要约 $2ndd_f$ 次乘加。批量大小和层数分别再乘 $B$ 和 $L$，词表投影还需另计。若把一次乘法和一次加法算作两个浮点操作，应将主要乘加项再乘二；报告 FLOPs 时必须说明口径。由此可见，当 $n$ 尚未远大于 $d$ 时，投影和 FFN 可能占据大量计算，不能只盯着二次注意力项。

朴素实现中，分数或概率张量形状为 $(B,h,n,n)$，单个张量含 $Bhn^2$ 个元素。取 $B=1$、$h=32$、$n=4096$，元素数为 $536870912$；若每个元素两字节，仅这一个张量就占 $1$ GiB，其中 $1\,\mathrm{GiB}=2^{30}$ 字节。训练还要保存其他激活和梯度，实际峰值取决于是否融合、重算与复用缓冲，不能从这个数字直接得出整个训练任务的显存。固定头数时常写注意力矩阵存储为 $O(n^2)$，若比较头数变化则应保留 $h$。

\subsection{键值缓存为什么成立}

\keyterm{键值缓存（key--value cache, KV cache）}保存每一层历史位置的键和值。它成立的核心是因果性：固定参数、位置规则与确定性推理时，追加未来词元不会改变已有位置可读到的内容。逐层归纳可知，旧位置每层的表示及其键值都不变。新词元只需计算自己的查询、键和值，并让自己的查询读取历史缓存以及当前键值。这个论证依赖因果结构；双向编码器追加词元后旧位置可能读取新内容，不能直接照搬。

通常不缓存历史查询，是因为接下来的输出只需要新位置的查询；历史查询所对应的历史输出已经计算完毕。缓存查询并非数学上禁止，而是一般没有必要。缓存也并不消除新查询与所有历史键的匹配：若当前历史长度为 $t$，单步注意力仍需 $O(td)$ 运算，另有当前词元的投影和 FFN。生成 $T$ 个词元、忽略初始提示长度时，缓存使投影与 FFN 从反复重算前缀的 $O(T^2d^2)$ 降至 $O(Td^2)$，密集注意力累计仍为 $O(T^2d)$。若完全每步重算整个前缀，则其注意力累计为 $O(T^3d)$；这里比较的是明确的两种实现，不能把它简化成整个生成过程由二次降为一次。

\keyterm{预填充（prefill）}同时处理已知提示的多个位置并建立缓存；\keyterm{逐步解码（decode）}处理后续新位置并扩展缓存。长提示预填充常能充分利用大矩阵运算，少量请求的逐步解码则可能更受权重与缓存读取带宽限制；这是常见性能情形，不是所有硬件和批量下的绝对分类。首词元延迟包含排队、预填充和首次输出等开销，后续每词元延迟反映解码阶段；吞吐量还受到并发和调度影响，应分开测量。

\subsection{多查询与分组查询}

令键值头数为 $h_{\rm kv}$，且每个头的键和值维度均为 $d_k$。\keyterm{多查询注意力（multi-query attention, MQA）}保留多个查询头，却共享一个键头和一个值头；\keyterm{分组查询注意力（grouped-query attention, GQA）}让每组查询头共享一对键值头。MHA 对应 $h_{\rm kv}=h$，MQA 对应 $h_{\rm kv}=1$，GQA 在二者之间提供配置选择\cite{ainslie2023}。常规分组要求 $h$ 能被 $h_{\rm kv}$ 整除。共享键值不意味着共享查询，各查询头仍能产生不同的权重分布。

忽略内存管理开销，缓存长度为 $n$、每个数值占 $s$ 字节时，全部层的 KV 缓存约为
\begin{equation}
\operatorname{Bytes}_{\rm KV}=2BLnh_{\rm kv}d_ks.
\label{eq:tr-kv-memory}
\end{equation}
因子二来自键和值，而不是前向和反向。若 $B=1$、$L=32$、$n=4096$、$h=32$、$d_k=128$、$s=2$，MHA 缓存为 $2$ GiB；取 $h_{\rm kv}=8$ 则约为 $0.5$ GiB。它不包含模型权重、临时缓冲、内存分块浪费或其他请求。KV 头数减少会降低缓存容量及理想情况下的读取量，却不能保证相同倍数的端到端加速，也不保证质量完全不变。

\subsection{精确注意力的执行优化与缓存管理}

FlashAttention 通过分块、片上复用和在线归一化，避免把完整注意力矩阵写入较慢的显存，再读回用于后续计算\cite{dao2022}。其核心是改变执行顺序和数据移动方式，而非默认删除远距离关系。它在数学上计算标准稠密注意力，浮点舍入顺序可能与朴素实现不同，因此“精确”不表示逐位相等。固定头维度时，它仍需处理二次数量级的查询--键关系，但可以显著降低中间显存和访存成本。

在线 Softmax 的合并公式可以直接验证。对同一查询的一组已处理键，保存最大分数 $m$、指数和 $\ell=\sum_j e^{s_j-m}$ 及未归一化加权向量 $u=\sum_j e^{s_j-m}v_j$。另一个键块给出 $(m_b,\ell_b,u_b)$，令 $m'=\max(m,m_b)$，则
\begin{equation}
\ell'=e^{m-m'}\ell+e^{m_b-m'}\ell_b,\qquad
u'=e^{m-m'}u+e^{m_b-m'}u_b,
\qquad o=u'/\ell'.
\label{eq:tr-online-softmax}
\end{equation}
这是把两个块统一到同一个指数基准后相加，因此不需要保存所有权重。公式描述数学核心，完整高性能实现还涉及分块尺寸、线程调度、反向重算和硬件约束。

PagedAttention 关注服务端 KV 缓存的分块存储与地址管理\cite{kwon2023}，减少连续预分配造成的浪费，并支持适当条件下的共享。它与 FlashAttention 所解决的问题不同，两者可以共同使用。连续批处理动态加入新请求并移除完成请求，提高设备利用率；前缀缓存复用相同前缀的计算，需要模型权重、适配器、词元、位置规则等相关状态兼容，不能只比较原始字符串。服务优化应同时报告首词元延迟、后续延迟、吞吐及峰值显存，因为单个指标改善不代表所有用户体验同时改善。

\section{反向传播、参考实现与验证方法}
\label{sec:tr-implementation}

\subsection{注意力梯度的可核查推导}

下面对有限分数且每行至少有一个有效键的情况推导，掩码固定且禁止位置不参与求导。设 $O=AV$，上游梯度 $G=\partial\mathcal L/\partial O\in\mathbb R^{n\times d_v}$。由矩阵乘法的微分，
\[
\frac{\partial\mathcal L}{\partial V}=A^{\mathsf T}G,
\qquad U_A=\frac{\partial\mathcal L}{\partial A}=GV^{\mathsf T}.
\]
将式\eqref{eq:tr-softmax-grad}逐行应用，令 $c_i=\sum_j A_{ij}(U_A)_{ij}$，则分数梯度 $D$ 满足
\begin{equation}
D_{ij}=A_{ij}\bigl((U_A)_{ij}-c_i\bigr),\qquad
\frac{\partial\mathcal L}{\partial Q}=\frac{DK}{\sqrt{d_k}},\qquad
\frac{\partial\mathcal L}{\partial K}=\frac{D^{\mathsf T}Q}{\sqrt{d_k}}.
\label{eq:tr-attention-backward}
\end{equation}
各梯度的形状分别与对应输入一致，可用作维度自检。由于 $\sum_jD_{ij}=0$，整体平移某一行分数不影响输出这一不变性也在梯度中得到体现。继续对 $Q=XW_Q$ 求导可得查询路径对 $X$ 的贡献为 $(\partial\mathcal L/\partial Q)W_Q^{\mathsf T}$；自注意力中还要加上键路径与值路径的贡献，不能只沿查询路径反传。

\subsection{单头注意力与增量掩码的参考代码}

下列 PyTorch 代码直接实现式\eqref{eq:tr-attention}，输入已经按头投影，统一形状为 $(B,h,n,d_k)$、$(B,h,m,d_k)$ 和 $(B,h,m,d_v)$。它使用 FP32 计算分数和归一化，并要求每行有有效键；目的在于清楚展示运算与掩码语义，不是替代融合内核。此处布尔值真表示允许读取，代码不含随机失活、不自动添加位置编码，也不自动构造填充掩码。若所有输入均为有限值，输出返回值张量的数据类型。

\begin{verbatim}
import math
import torch

def attention_reference(q, k, v, allowed=None):
    # q: [B, H, N, Dk]; k: [B, H, M, Dk]
    # v: [B, H, M, Dv]; allowed: broadcastable to [B,H,N,M]
    assert q.ndim == k.ndim == v.ndim == 4
    assert q.shape[:2] == k.shape[:2] == v.shape[:2]
    assert q.shape[-1] == k.shape[-1]
    assert k.shape[-2] == v.shape[-2]
    score = q.float() @ k.float().transpose(-2, -1)
    score = score / math.sqrt(q.shape[-1])
    if allowed is not None:
        allowed = torch.broadcast_to(allowed.bool(), score.shape)
        if not bool(allowed.any(dim=-1).all()):
            raise ValueError("Each query needs a visible key")
        score = score.masked_fill(~allowed, float("-inf"))
    prob = torch.softmax(score, dim=-1)
    return (prob @ v.float()).to(v.dtype)

def incremental_causal_mask(past_len, new_len, device):
    query_pos = past_len + torch.arange(new_len, device=device)
    key_pos = torch.arange(past_len + new_len, device=device)
    return key_pos[None, :] <= query_pos[:, None]
\end{verbatim}

例如已有三个缓存位置、一次输入两个新位置时，返回掩码的两行为 $(1,1,1,1,0)$ 和 $(1,1,1,1,1)$。单词元增量输入时，已有缓存加当前词元都应可见。上面的函数假定查询与键值有相同头数；GQA 应依据组映射匹配查询头和键值头，不能仅将头数不同的张量直接相乘。左填充、样本拼接、滑动窗口或不同请求各自拥有不同缓存长度时，还需要结合每个样本的有效位置构造掩码。

\subsection{从现象到错误位置的验证路线}

实现验证首先用短序列和小维度，固定参数并关闭随机失活。手算例子检查 Softmax 轴和掩码，改变未来输入后比较较早位置的输出，检查因果性；把完整前向最后位置的逻辑值与逐词元缓存前向对比，检查缓存与位置偏移；把独立样本与填充批处理结果比较，检查填充和有效位置选择。不同计算路径存在浮点误差，应报告最大绝对或相对误差及精度，而不是要求不同融合实现逐位相同。

梯度检查在小规模、较高精度、无随机失活和无不可微分边界时，可比较自动微分、式\eqref{eq:tr-attention-backward}和有限差分。有限差分步长过大引入截断误差，过小放大舍入误差，因此需要观察一段合理步长范围。测试通过只能支持被覆盖的形状和配置；从单头无缓存通过，不能跳到多头、混合精度、变长批处理都正确的结论。性能测试应在正确性之后，并区分预热、同步、数据传输和真实计算时间。

训练不收敛时应先确认数据与标签，再检查数学与数值。固定极小训练集能否明显降低损失是有用诊断，但不是泛化能力测试。如果训练损失异常接近零，要检查是否预测了当前输入词元、是否移位两次、未来键是否泄漏，以及验证集是否重复。若出现非有限值，应沿层记录激活尺度、分数极值、归一化分母和梯度范数，定位首次出现问题的运算；盲目换优化器可能暂时掩盖错误，却不能建立可信的训练过程。

\section{面试问答：基础概念与信息流}
\label{sec:tr-qa-basic}

\begin{enumerate}[label=\arabic*.,series=trquestions]
\item \textbf{怎样用两分钟说明 Transformer 的工作原理？}
\textbf{答：}先把离散词元映射为向量，并引入位置机制；每层用注意力决定各位置读取哪些内容，用前馈网络变换当前位置的特征，再由残差与归一化组织深层计算。对于因果语言模型，未来位置被屏蔽，最后的表示投影到词表形成下一词元分布，训练最小化真实下一词元的负对数似然，生成则反复选择新词元并追加。追问时应能指出 $QK^{\mathsf T}$ 负责匹配、$AV$ 负责汇总，以及模型架构、训练目标和采样策略属于不同层次。

\item \textbf{Transformer 相比循环神经网络有什么优势和代价？}
\textbf{答：}循环神经网络（recurrent neural network, RNN）按时间递推隐藏状态，训练中同一序列的位置之间存在状态依赖；自注意力可以在同一层并行处理已知位置，远处位置之间也可直接交换信息。代价是标准稠密注意力处理所有位置对，长序列计算与存储较重。RNN 可用固定大小状态流式更新，而标准带完整 KV 缓存的 Transformer 状态随长度增长。因此应结合训练吞吐、流式限制、任务长度与硬件讨论，不能回答成 Transformer 在任何场景都更快、更准确。

\item \textbf{Transformer 与卷积网络的归纳偏置有什么不同？}
\textbf{答：}卷积神经网络（convolutional neural network, CNN）显式规定局部邻域与共享卷积核，适合表达局部结构；标准注意力的连接权重随输入内容改变，单层即可覆盖整个允许序列。全局可达不等于一定学到有效远程依赖，局部先验也不等于表达能力不足。窗口注意力、卷积与注意力混合等结构正是在计算预算和任务先验之间做选择。比较时应控制数据量、参数量和训练预算，不能仅用网络名称解释实验差异。

\item \textbf{为什么查询、键和值要做三个投影？可以共用吗？}
\textbf{答：}三个投影把匹配条件、可被匹配特征和返回内容分开，使模型能够学习有方向的关系。它们可以按设计部分共享，甚至直接使用同一个输入，但会改变参数化与表达限制。特别是 $Q=K$ 时原始点积分数具有对称性，独立投影则不受这一限制；输出权重是否对称还取决于逐行归一化。面试中不能把“三个投影”解释成三份不同的训练数据，它们通常来自同一个隐藏表示。

\item \textbf{为什么称为自注意力？它与交叉注意力的本质差别是什么？}
\textbf{答：}“自”描述查询和键值的来源相同，不表示一个词元只关注自己。交叉注意力中查询来自接收信息的一侧，键和值来自提供信息的一侧，输出长度由查询侧决定。例如目标序列长 $n$、源序列长 $m$，分数矩阵就是 $n\times m$，输出仍有 $n$ 行。两种机制可用同一个注意力公式实现，只是投影输入、长度和可见范围不同；交叉注意力不是把两个句子的词元逐个一一配对。

\item \textbf{注意力是不是一种概率模型？权重能否表示置信度？}
\textbf{答：}Softmax 权重在没有随机失活时构成每个查询上的离散概率向量，但这首先是信息混合规则，并未自动规定某个真实随机事件。最大注意力权重不是答案正确概率，注意力熵也不是可靠性评估的通用替代。若要讨论置信度，需要明确预测事件、观察标签和校准方式。一个模型可能高度集中地读取错误位置，也可能在分散读取多段证据后给出正确答案。

\item \textbf{注意力层中的值向量为什么不能被查询和键完全替代？}
\textbf{答：}查询和键决定“读多少”，值决定“读出什么”。只保留分数矩阵会得到位置关系，却没有规定向下一层传递的内容表示。可以设计让值等于键或原始输入的结构，但那是把两个角色绑定，并没有取消内容汇总这一步。数学上 $d_k$ 与 $d_v$ 可不同，$QK^{\mathsf T}$ 只要求查询和键宽度一致，最终输出宽度由值及输出投影决定。

\item \textbf{自注意力是否完全取代了特征提取？}
\textbf{答：}没有。词元嵌入、查询键值投影、前馈网络和输出头都在构造或变换特征，注意力也通过输入相关权重形成新表示。把它简单称为“只做加权平均”遗漏了值的可学习投影和多层非线性组合。单头单次输出的凸组合性质只是局部代数性质，不足以限制整个带残差与前馈网络的函数族，更不能据此认定它只能复制训练样本。
\end{enumerate}

\section{面试问答：公式、维度与数学性质}
\label{sec:tr-qa-math}

\begin{enumerate}[resume*=trquestions]
\item \textbf{为什么缩放因子是 $\sqrt{d_k}$，而不是 $\sqrt d$ 或序列长度？}
\textbf{答：}每个分数是一个头内部 $d_k$ 项乘积之和，尺度分析针对的是这些求和项，因此应使用单头维度。独立、零均值、单位方差假设下，未缩放点积方差为 $d_k$，除以平方根后方差为一。模型宽度 $d$ 只有在单头维度等于它时才可直接使用；序列长度影响候选数和 Softmax 分配，却不是这段方差推导里的求和维度。真实训练中独立假设通常不严格成立，所以这是设计动机而非分布定律。

\item \textbf{Softmax 应该沿哪个维度计算？沿错维度有什么后果？}
\textbf{答：}对 $(B,h,n,m)$ 的分数沿最后的键维度 $m$ 计算，使每个查询在允许键之间分配总权重一。若沿查询维归一化，表达的将是每个键在不同查询间如何分配权重，改变了运算含义，也破坏通常的逐查询凸组合解释。最直接检查是构造两个查询、三个键的非对称例子，验证每行和为一及手算输出；只检查输出形状无法发现这种错误。

\item \textbf{自注意力的权重矩阵一定对称吗？}
\textbf{答：}不一定。独立的查询与键投影使 $QK^{\mathsf T}$ 就可能不对称；即使令 $Q=K$，不同分母的逐行 Softmax 也通常破坏对称性。反例是分数矩阵 $\left(\begin{smallmatrix}0&0\\0&1\end{smallmatrix}\right)$，归一化后 $A_{12}=1/2$，而 $A_{21}=1/(1+e)$。因果掩码还会引入方向性，因此不能把注意力矩阵当作无向图邻接矩阵而不检查条件。

\item \textbf{多头注意力是否等价于一个宽度相同的大头？}
\textbf{答：}通常不等价。多头有多套独立的 Softmax 权重，不同内容通道可以从不同位置汇总信息；一个大头只有一套位置分配，再宽的值向量也共用它。特定参数或输入下，两者可能碰巧给出相同输出，但这不能证明函数族等价。若只把输入机械切片、没有可学习投影，也只是某种受限实现，不能代表完整多头机制。

\item \textbf{固定模型宽度时，增加头数会增加多少参数？}
\textbf{答：}在标准 $d_k=d_v=d/h$ 配置下，合并后的四个投影仍约有 $4d^2$ 个参数，头数增加不使这一主项成倍增长。每头宽度会相应减少，朴素保存的权重张量却含 $hn^2$ 个元素，因此中间显存和内核效率可能变化。若保持单头宽度不变而增加头数，总投影宽度就增长，结论不同。回答必须先固定比较中的不变量。

\item \textbf{给定 $B=2,n=128,d=768,h=12$，各张量形状是什么？}
\textbf{答：}单头宽度为 $64$。合并投影后的查询、键和值都是 $(2,128,768)$，拆头并交换轴后是 $(2,12,128,64)$，分数与权重形状为 $(2,12,128,128)$，乘值后的输出为 $(2,12,128,64)$，合并后回到 $(2,128,768)$。这些结论假设自注意力和标准等宽头；交叉注意力的键长度应改为源长度，GQA 的键值头数还要单独改变。

\item \textbf{如何推导注意力对查询、键和值的梯度？}
\textbf{答：}先把计算拆成 $O=AV$、$A=\operatorname{softmax}(S)$、$S=QK^{\mathsf T}/\sqrt{d_k}$。对矩阵乘法求导得到 $\partial\mathcal L/\partial V=A^{\mathsf T}G$ 与 $\partial\mathcal L/\partial A=GV^{\mathsf T}$，再逐行乘 Softmax 雅可比得到分数梯度，最后分别右乘 $K$ 或转置后乘 $Q$，详见式\eqref{eq:tr-attention-backward}。自注意力的输入同时流向三个投影，三条梯度必须相加；只背最后公式而不说明形状容易把转置写反。

\item \textbf{注意力输出为什么不会总是趋于所有词元的平均？}
\textbf{答：}均匀平均只在所有可见分数相同或十分接近时出现。查询键投影、内容差异、位置项和掩码均可产生不同分数，训练也会通过任务梯度改变匹配。即使某个头暂时接近平均，其值投影与其他头仍可能有用。反过来，权重非常尖锐也不自动代表质量更好；是否适当取决于任务需要整合多个证据还是选择少数位置。
\end{enumerate}

\section{面试问答：位置、归一化与架构设计}
\label{sec:tr-qa-architecture}

\begin{enumerate}[resume*=trquestions]
\item \textbf{没有位置编码时，Transformer 能否区分词序？}
\textbf{答：}对于没有位置相关掩码、所有位置共享参数的标准自注意力堆叠，它具有式\eqref{eq:tr-permutation}的置换等变性，不能独立赋予任意排列不同的顺序结构。但因果掩码已经改变可见范围，使这一前提不成立，不能把结论推广成所有无显式位置编码模型都不可能表达任何顺序。面试重点是说出等变性及其假设，并区分词元级输出随排列变化和池化后整体输出不变。

\item \textbf{位置嵌入为什么常与词元嵌入相加，而不是拼接？}
\textbf{答：}相加保持宽度 $d$，便于沿用残差与投影形状；拼接也完全可以，但会增加宽度或需要额外投影。相加并不保证模型总能无损分离内容与位置，它是一种让训练共同组织表示的设计。采用拼接是否更好需要控制参数与计算预算后比较，不能从“信息多一份”直接推出优势。RoPE 则采用另一条路径，在匹配向量上进行旋转，而非必须在输入嵌入处做加法。

\item \textbf{RoPE 为什么具有相对位置性质？}
\textbf{答：}关键恒等式是 $R(i\omega)^{\mathsf T}R(j\omega)=R((j-i)\omega)$，使旋转后查询键点积中的显式位置因素变成差值。应补充两个限定：分数仍依赖内容向量，并不是距离函数；“相对位置形式”也不证明任意长度外推可靠。若面试进一步追问长上下文，需要讨论训练距离分布、频率缩放、位置索引与评测，而不是只重复“旋转不改变范数”。

\item \textbf{为什么通常使用 LayerNorm，而不是 BatchNorm？}
\textbf{答：}批归一化（batch normalization, BatchNorm）依赖选定批量轴及其他归约轴的统计，变长序列、不同批量与流式推理会使这些统计的管理更复杂。标准 LayerNorm 对每个词元的特征维独立计算，训练与推理不依赖运行中的批统计量，较适合这种表示组织。若错误地跨时间计算归一化统计，还可能引入未来信息。不能说 BatchNorm 在序列模型中数学上不可用，真正问题是统计轴、掩码和部署方式是否一致。

\item \textbf{Pre-LN 一定优于 Post-LN 吗？}
\textbf{答：}没有这种无条件保证。Pre-LN 提供更直接的残差梯度路径，在一些设置中更容易稳定训练；Post-LN 对每次残差相加后的表示实施归一化，梯度与表示尺度的演化不同。最终结果还受深度、初始化、学习率与残差缩放影响。合理回答是说明公式和稳定性机制，再指出文献\cite{xiong2020}的结论有具体理论与实验前提，不能直接宣称所有 Pre-LN 都可取消预热。

\item \textbf{RMSNorm 相比 LayerNorm 少了什么？}
\textbf{答：}它不减均值，而是用均方根缩放；常见形式也没有平移参数。LayerNorm 对所有坐标加同一常数具有标准化层面的平移不变性，RMSNorm 一般没有这一性质。二者都在特征维控制尺度，但归一化后的几何含义不同。是否速度更快需要测量，因为归约、融合、内存访问与硬件会影响收益；不能只数公式里少了一个减法就预测完整模型加速比例。

\item \textbf{FFN 为什么先升维再降维？一定要四倍宽吗？}
\textbf{答：}中间层提供更多非线性通道，之后投影回模型宽度以便残差相加。四倍宽是常见配置，不是数学必需；门控 FFN 为匹配参数预算还可能采用不同宽度。提升宽度、增加层数和改变激活函数各自改变表达与计算，必须按预算讨论。单纯升维而无非线性仍可合并为一个线性变换，无法获得通常希望的非线性表达增益。

\item \textbf{仅解码器模型为什么也能做分类、翻译和摘要？}
\textbf{答：}这些任务可以写成在输入条件下生成目标文本，因果模型先读取序列化输入，再生成标签或答案；也可以从某个表示接分类头。它能完成这些任务，并不意味着编码器或编码器--解码器失去价值，后者可在双向表示、源目标长度差异或特定推理预算下有不同取舍。比较应固定任务协议与资源，不能把架构流行程度当作严格优越性的证明。
\end{enumerate}

\section{面试问答：训练目标与数值稳定性}
\label{sec:tr-qa-training}

\begin{enumerate}[resume*=trquestions]
\item \textbf{训练时为什么能并行，生成时却通常逐词元进行？}
\textbf{答：}教师强制提供全部真实前缀，每个位置的输入已经知道，因此同一层可以同时计算所有位置；因果掩码限制读取关系。生成时下一步输入包含刚选出的新词元，这个离散值在前一步结束前未知，形成序列依赖。并行计算所有位置的逻辑值不等于一次就能得到任意长的未知续写。投机解码等方法可改变执行组织，但需要额外验证或其他假设，不能据此抹去自回归条件依赖。

\item \textbf{因果掩码为什么通常保留对角线？}
\textbf{答：}位置 $i$ 输入的是当前已知词元 $x_i$，其输出用于预测 $x_{i+1}$，所以读到自己是合法条件。若同位置输出直接预测同一个输入词元且没有遮挡或移位，就可能成为复制任务。用一个带起始和结束符的短序列写出输入与标签，通常比背“上三角置负无穷”更能证明理解。掩码朝向还依赖行表示查询、列表示键这一约定，不能脱离轴定义讨论。

\item \textbf{为什么只对回答计算损失，却仍要让模型读取提示词？}
\textbf{答：}损失掩码规定哪些输出接受监督，注意力掩码规定哪些输入可以提供条件。回答的条件概率需要依赖问题与上下文，所以提示词应参与前向信息流；不对提示词位置计分，只表示本轮不直接优化这些位置的预测目标。提示词的嵌入与中间表示仍可能通过回答位置的损失收到梯度，不能把忽略标签误解为整段输入完全不参与训练。

\item \textbf{交叉熵下降是否说明生成能力一定提高？}
\textbf{答：}它说明在指定数据和计分协议下，真实目标获得更高的平均对数概率；不直接证明任务正确率、事实可靠性或长答案质量同步提高。训练损失还可能因过拟合、数据重复或泄漏下降。应同时检查验证似然、任务指标、生成设置与错误样本。不同数据配比或只计算回答损失也会改变平均值的含义，不能跨协议只比较一个损失数字。

\item \textbf{梯度累积什么时候等价于更大的批量？}
\textbf{答：}在相同参数上累积各微批次正确归一化的梯度，并在全部累积完成后只更新一次时，它对应合并目标的梯度；还需考虑随机失活、依赖批统计的层和浮点求和次序。有效词元数不同要按词元加权，梯度裁剪通常应在全部累积完成后进行，否则会改变方向与大小。学习率调度和优化器步数也必须以实际更新为准，而不能把每个微批次都当作一次参数更新。

\item \textbf{损失变成 NaN，应优先检查什么？}
\textbf{答：}先检查输入和掩码是否产生空有效行、损失分母是否为零、逻辑值是否已有无穷，再追踪低精度溢出、归一化分母、学习率与梯度。保存首次异常的批次及前一个可恢复状态，缩小到可复现例子，通常比直接降低学习率更有效。某些融合内核对全屏蔽行有专门处理，也不能因此忽略语义检查，因为更换后端可能暴露不同表现。

\item \textbf{为什么微调后训练集很好，验证集却变差？}
\textbf{答：}可能是过拟合、训练验证分布不一致、训练数据质量差、提示模板变化或通用能力遗忘。应比较微调前后在同一冻结评测协议下的结果，核对模板、分词、截断与解码设置，进一步按领域和长度分组。更小学习率、早停或数据混合是候选干预，不是诊断结论；要通过对照确定问题来源，避免将格式不兼容误判为模型能力退化。

\item \textbf{推理模式和关闭梯度有什么区别？}
\textbf{答：}模块的评估模式通常改变随机失活等层的行为，关闭自动微分则减少计算图与梯度相关开销，两者不是同一个开关。只关闭梯度，随机失活仍可能存在；只切评估模式，也可能仍构建计算图。函数式注意力接口还可能直接按传入的失活概率执行，需在评估时显式传零。要做缓存一致性测试，应同时控制随机性、梯度记录和模型参数状态。
\end{enumerate}

\section{面试问答：生成、缓存与服务行为}
\label{sec:tr-qa-cache}

\begin{enumerate}[resume*=trquestions]
\item \textbf{KV 缓存究竟缓存了什么，为什么每一层都需要？}
\textbf{答：}缓存每一层已经处理位置的键和值，而不是只保存词元嵌入或最后一层隐藏状态。各层输入表示不同，投影矩阵也不同，新查询必须读取同一层对应的历史键值，所以通常每层有自己的缓存。启用 RoPE 时常缓存已经完成位置旋转的键，具体约定必须与读取过程一致，不能重复旋转。缓存有效性来自固定参数、确定性计算与因果可见范围，追加未来词元不改变历史表示。

\item \textbf{为什么不缓存查询？使用缓存后是否不用计算注意力？}
\textbf{答：}后续步骤不再需要历史位置的输出，因而通常不需要历史查询；当前新查询仍要与所有可见历史键比较，并汇总对应的值。缓存主要避免历史位置的投影、前馈和注意力结果重复计算，不会消除当前查询读取历史的工作。因此上下文变长时单步解码仍可能变慢，显存和带宽需求也会上升。若某个系统保存额外状态用于调度或诊断，也不改变这个基本依赖分析。

\item \textbf{有缓存时，整个生成过程的复杂度是线性吗？}
\textbf{答：}标准全上下文注意力不是。固定宽度下，第 $t$ 步需读取约 $t$ 个键值，累计 $\sum_{t=1}^{T}t=O(T^2)$；新词元自身的投影与 FFN 累计才是对 $T$ 线性。若有长度为 $P$ 的初始提示，还会出现约 $PT$ 的生成期注意力项及独立预填充代价。只有在固定窗口、压缩状态或其他改变访问模式的假设下，才可能得到不同的增长关系。

\item \textbf{缓存大小如何随批量、上下文长度和头数变化？}
\textbf{答：}使用式\eqref{eq:tr-kv-memory}逐项代入，容量对批量、层数、长度、键值头数、单头宽度和每元素字节数都线性增长。在固定 $d$ 的 MHA 中，$hd_k=d$，单独改变头数并保持总宽度不变，并不会改变主要 KV 元素数；GQA 减少的是键值总宽度。真实服务还需区分所有请求的有效长度与分配容量，分块浪费和共享前缀会改变实际占用。

\item \textbf{为什么长提示首词元很慢，但之后每个词元较快？}
\textbf{答：}首次要完成整个提示的多层预填充并建立缓存，后续通常只计算新位置。不过随着历史增长，后续步骤读取的缓存也在增长，不能预期每步始终一样快。若首词元变慢，还应区分排队、模型加载、编译预热、数据传输与预填充本身；若后续变慢，应关注并发、缓存带宽和批处理调度。只报告整段生成总时间会把这些因素混在一起。

\item \textbf{贪心解码为什么不保证得到最高概率的完整句子？}
\textbf{答：}完整序列概率是条件概率的乘积，第一步最优选择可能导致后续概率很低。例如第一步 $p(a)=0.6$、$p(b)=0.4$，若在 $a$ 后第二步最高概率为 $0.5$，在 $b$ 后某个词元概率为 $0.9$，两步最佳概率分别为 $0.30$ 和 $0.36$，贪心会错过后者。这个反例固定生成长度为两步；变长生成还要处理结束符和长度评分。束搜索能扩大探索，却仍受束宽限制。

\item \textbf{温度、Top-$k$ 和 Top-$p$ 应怎样区分？}
\textbf{答：}温度改变所有候选的相对概率比例，Top-$k$ 固定保留候选数，Top-$p$ 按累计概率选择可变大小的集合。若同时使用，执行顺序会影响最终分布，例如先调温再按累计概率截断可能改变保留集合。低温不意味着答案一定正确，高温也不等于真正的知识创新。实验比较应固定随机种子、采样次数与长度限制，并报告任务结果的波动。

\item \textbf{为什么改了一点提示词，原有前缀缓存可能失效？}
\textbf{答：}因果网络中，一个位置的变化可以影响其后所有位置的表示，因此只有变化之前保持一致且相关状态兼容的前缀能直接复用。字符串表面相似还不够，分词、模板、特殊词元和位置索引都应相同；模型权重或适配器变化也可能使原缓存不适用。正确复用单位是兼容计算所产生的公共前缀状态，而不是两个请求在自然语言意义上“差不多”。
\end{enumerate}

\section{面试问答：效率与复杂度比较}
\label{sec:tr-qa-efficiency}

\begin{enumerate}[resume*=trquestions]
\item \textbf{一句话说注意力是 $O(n^2)$，遗漏了哪些信息？}
\textbf{答：}它遗漏头宽度、批量、层数，以及投影和前馈网络。标准自注意力核心乘法是 $O(n^2d)$，线性投影是 $O(nd^2)$，FFN 是 $O(ndd_f)$；朴素概率矩阵存储是 $O(hn^2)$。还要区分训练整段前向、预填充和单步解码，因为查询长度不同。完整答案应先定义比较变量，再说明瓶颈可能由哪个项主导。

\item \textbf{FlashAttention 是否把二次注意力变成了线性注意力？}
\textbf{答：}没有。它优化标准稠密注意力的数据移动与中间存储，通常仍处理二次数量级的位置关系。中间显存随长度的增长方式改善，不等于算术运算阶数一起下降。线性注意力则通常改变核函数、归一化或可交换的计算结构，属于另一类设计。回答时还应区分数学等价与浮点逐位一致，避免把数值舍入差异误判为算法采用了近似稀疏化。

\item \textbf{FlashAttention 和 PagedAttention 的区别是什么？}
\textbf{答：}前者主要提高注意力算子执行效率，通过分块和在线 Softmax 降低中间矩阵的显存读写；后者主要管理服务中的 KV 缓存布局，减少存储浪费并支持相应的共享与调度。二者可以互补，不能仅根据名字中都有 Attention 就当成同一种算法。选型还要考虑后端是否兼容、请求长度分布与并发，不能脱离系统环境保证固定加速倍数。

\item \textbf{GQA 相比 MHA 的主要收益和代价是什么？}
\textbf{答：}多个查询头共享键值投影，减少键值参数、缓存容量和潜在读取流量，因此常有利于解码部署。代价是键值表示的独立性受到约束，质量与训练适配需要评估；查询头数量不变并不意味着所有计算都不变，也不意味着计算量按缓存比例下降。实际推理若把共享键值物理复制多次，还可能失去部分存储收益，应区分逻辑广播与真实分配。

\item \textbf{模型参数少了，推理一定更快吗？}
\textbf{答：}不一定。速度还受算子形状、内存带宽、并行度、内核支持、通信与调度影响。一个参数更少但由许多小算子组成的模型，可能比矩阵形状更友好的模型慢；稀疏结构也需要高效执行才能兑现理论节省。应在相同硬件、批量、输入输出长度和精度下测量首词元与解码延迟，并同时报告质量，避免只用参数量预测端到端体验。

\item \textbf{为什么推理能放进显存，训练却放不下？}
\textbf{答：}训练除权重外还保存梯度、优化器状态和用于反向的激活，部分混合精度方案另有主权重副本。推理没有这些训练状态，却可能有很大的 KV 缓存。以全为四字节数值的普通 Adam 为例，权重、梯度、一阶矩和二阶矩合计约为每参数十六字节，尚未计激活；这只是特定核算口径，不可套用所有量化或分片方案。显存溢出应先区分状态占用和激活占用。

\item \textbf{长序列训练显存不足，应怎样决定优化顺序？}
\textbf{答：}先测量权重、优化器状态、激活和临时张量的占比。注意力矩阵主导时考虑融合精确注意力，激活主导时考虑检查点与较小微批量，训练状态主导时考虑状态分片或参数高效适配。截断长度、稀疏窗口和压缩上下文会改变任务或信息范围，需要单独评估质量。不能把所有方案当成等价的“省显存开关”，因为它们分别交换计算、通信、有效批量或模型行为。

\item \textbf{吞吐提高是否意味着单个用户等待时间更短？}
\textbf{答：}不一定。较大批量可能增加单位时间处理的词元数，却让请求等待凑批或竞争计算资源；持续批处理也可能在不同请求间分配不均。应该报告请求到达模式、并发量、输入输出长度分布以及延迟分位数，区分服务整体吞吐与单请求首词元、每词元延迟。只有在同一负载和质量约束下比较，指标改善才有明确含义。
\end{enumerate}

\section{面试问答：实现、排错与实验设计}
\label{sec:tr-qa-debug}

\begin{enumerate}[resume*=trquestions]
\item \textbf{怎样证明自己的因果掩码没有泄漏未来信息？}
\textbf{答：}固定参数并关闭随机失活，构造两条前缀相同、后缀不同的等长序列，比较所有相同前缀位置的输出；正确因果实现应在浮点容差内保持一致。还可直接检查小尺寸允许矩阵和未来位置对较早输出的梯度是否为零。测试应覆盖多层，因为某一层的错误也可能经后续传播；同时注意跨位置归一化或数据预处理也可能造成泄漏，不能只检查注意力那一行代码。

\item \textbf{启用 KV 缓存后输出不一致，应按什么顺序排查？}
\textbf{答：}先关闭采样与随机失活，比较逻辑值而非最终文本；然后检查输入是否重复包含已缓存前缀、每层缓存长度、RoPE 绝对位置、增量因果偏移和填充规则。再检查头维重排、GQA 分组与缓存数据类型。逐层比较可以找到首次偏离的位置。最终词元相同不表示逻辑值完全正确，最终文本不同也可能只是微小数值差异被采样放大，因此诊断对象必须选对。

\item \textbf{布尔掩码的真值到底代表允许还是禁止？}
\textbf{答：}取决于接口。本文参考代码中真表示允许。PyTorch 官方文档明确区分：函数式缩放点积注意力的布尔掩码真表示参与，而多头注意力模块的键填充掩码真表示需要屏蔽\cite{pytorchsdpa}。不要依靠变量名推断语义，迁移实现时应重新核对版本文档。最小验证是只允许一个键，此时输出应等于那个值向量；这比观察大模型生成是否“看起来正常”更直接。

\item \textbf{为什么左填充和右填充可能影响批量生成结果？}
\textbf{答：}如果无条件取张量最后一列作为下一词元分布，右填充的短样本最后一列可能是占位位置；左填充则需要正确处理位置编号和被屏蔽的查询行。两种方式在合适的掩码、位置与索引实现下都可组织正确计算，问题不在于某一种排列天然不合法。应比较每个样本独立运行和批量运行的有效位置输出，同时确认结束符没有被错误当作普通填充参与预测。

\item \textbf{转置以后直接重塑张量，为什么容易出错？}
\textbf{答：}张量的逻辑轴顺序和底层连续存储不是一回事，转置常改变步幅而不移动数据。某些重塑接口要求连续布局，另一些可能复制数据，但无论是否报错，都必须保证最终元素排列符合“同一位置拼接各头”的定义。应在小张量中给每个元素不同编号，手工追踪拆头和合头前后的对应关系。只看维度乘积相同，不能证明语义顺序正确。

\item \textbf{如何判断加速内核结果正确？}
\textbf{答：}先与清楚的高精度参考实现比较前向结果，再按需要比较反向梯度，覆盖因果、非因果、交叉长度、不同头数及边界掩码。记录最大误差、数据类型和后端，不应对所有浮点实现要求逐位一致。若差异随长度或分数幅度异常放大，应进一步定位归一化、缩放与累积精度。正确性达标后才测性能，否则更快地计算错误结果没有实际价值。

\item \textbf{训练集能够被小模型记住，是否证明实现正确？}
\textbf{答：}它是排错信号，说明优化过程至少能降低某个目标，但不构成完整证明。错误的标签移位、未来泄漏或对填充符的过度预测也可能使损失很低。还需因果性、梯度、缓存等价性和样本独立性检查，再在无泄漏验证集上评估泛化。小数据过拟合测试应与结构性测试配合，而不能替代对计算图语义的理解。

\item \textbf{怎样设计“增加头数是否有用”的公平实验？}
\textbf{答：}先固定模型宽度、层数、训练数据与计算预算，明确头数变化时单头维度如何调整；保持优化与评测协议可比，记录参数、真实训练速度和显存。用多个随机种子或重复实验观察结果波动，并按任务与长度分析。若增加头数同时增加了总宽度，就不能把全部收益归因于头数。即使某个任务改善，也应把结论限制在测试配置，而不是推断头越多越好。
\end{enumerate}

\section{面试问答：扩展方法与综合判断}
\label{sec:tr-qa-advanced}

\begin{enumerate}[resume*=trquestions]
\item \textbf{LoRA 如何适配 Transformer，为什么能减少训练状态？}
\textbf{答：}低秩适配（low-rank adaptation, LoRA）冻结原矩阵 $W\in\mathbb R^{d_{\rm in}\times d_{\rm out}}$，学习 $\Delta W=cAB$，其中 $A\in\mathbb R^{d_{\rm in}\times r}$、$B\in\mathbb R^{r\times d_{\rm out}}$，$r$ 较小，$c$ 为约定的缩放因子。可训练参数从完整矩阵降为 $r(d_{\rm in}+d_{\rm out})$，相应梯度和优化器状态减少，仍需要基础权重和一定激活。目标模块可以是注意力投影或 FFN，效果需验证；低秩限制针对更新，不表示原矩阵必须低秩，方法来源见文献\cite{hu2021}。

\item \textbf{量化是否一定同时减少显存、提高速度并保持质量？}
\textbf{答：}量化（quantization）以较少位数近似表示权重、激活或缓存，通常可以减少相应张量的存储，但还需缩放因子、零点或其他元数据。速度收益取决于硬件是否支持对应低位运算、是否发生反量化及数据移动；质量变化取决于误差分布和敏感模块。权重量化不自动缩小 KV 缓存，缓存量化也不等于优化器状态量化，回答时必须指出被量化的对象。

\item \textbf{稀疏注意力与线性注意力分别改变了什么？}
\textbf{答：}稀疏注意力限制允许的位置对，如局部窗口、块连接或少量全局位置，以减少访问；它可能改变长程信息到达某位置所需的层数。线性注意力的一类方法用可分解核表示重排求和，例如把适当形式的计算写成 $\phi(Q)(\phi(K)^{\mathsf T}V)$，并处理相应归一化。标准逐行 Softmax 一般不能直接这样结合，因为非线性作用在完整分数上。两类方法的复杂度与质量都依赖具体结构，不能只凭名称判断。

\item \textbf{混合专家模型为什么要区分总参数与激活参数？}
\textbf{答：}混合专家（mixture of experts, MoE）通常让路由器为每个词元选择少数专家，专家可替换或扩展 FFN。模型保存的全部专家参数决定总容量与部分存储需求，每次真正参与某个词元计算的专家决定激活计算量。两者不同，部署还涉及路由不均、专家通信和容量限制。Switch Transformer 是这种条件计算的一种具体设计\cite{fedus2021}；专家并非一定对应清晰的人类知识领域，也不能由激活参数少直接推出单请求低延迟。

\item \textbf{支持很长上下文，是否意味着模型能有效利用其中全部信息？}
\textbf{答：}不是。长度上限首先可能只是位置表示和内存允许输入的范围；有效利用还受训练长度、证据位置、干扰信息和任务难度影响。评测应把证据放在不同位置，改变长度与干扰量，既测简单检索也测多段组合推理，并与短上下文或检索增强基线比较。一次准确找回关键词不证明具备同等长度的可靠推理，位置公式可外推也不等于模型行为能外推。

\item \textbf{能否直接平均所有隐藏状态作为句向量？}
\textbf{答：}可以构造这样的表示，但其相似度是否有用取决于训练目标、所选层、归一化和池化方式。填充位置必须排除；因果模型各位置可见上下文长度不同，双向编码器则有不同的信息条件。生成模型隐藏状态并未必经过句子相似度或检索目标优化，因此不能仅凭向量维度合适就宣称适合检索。应在独立匹配或检索任务上比较池化策略和专门训练的表示模型。

\item \textbf{Transformer 为什么能够处理图像和多模态输入？}
\textbf{答：}注意力接收的是向量序列，并不要求向量一定来自文字词表。视觉 Transformer（vision transformer, ViT）可把图像切成块并映射为向量，结合位置表示进入编码器\cite{dosovitskiy2020}；多模态模型还可把视觉或音频编码结果投影到语言模型空间，或通过交叉注意力读取。模态对齐、时空位置和训练目标依然必要，仅把不同向量拼接不保证模型自动理解它们的对应关系。

\item \textbf{面试中如何说明自己确实理解并实现过 Transformer？}
\textbf{答：}选择一个实际完成的最小项目，说明数据怎样分词和移位、各张量的形状、掩码允许谁读取谁、损失按什么单位平均，再展示因果性、梯度或缓存等价性检查，以及真实记录的训练和推理指标。若没有运行某项实验，应明确把它列为验证计划，不能编造提升比例。一个能定位错误、解释基线并界定结果范围的小项目，比罗列大量模型名更能展示可迁移能力。
\end{enumerate}

\section{综合计算题与复习方法}
\label{sec:tr-exercises}

\begin{example}[一层参数量与训练计算的核算]
考虑 $d=512$、$h=8$、$d_f=2048$ 的标准 MHA 加普通 FFN，仅计算矩阵参数，不含偏置和归一化。一层注意力参数为 $4\times512^2=1048576$，FFN 参数为 $2\times512\times2048=2097152$，合计 $3145728$。如果一条序列长 $n=1024$，注意力两个核心乘法约为 $2n^2d=1073741824$ 次乘加，四个投影也约为 $4nd^2=1073741824$ 次乘加，而 FFN 约为 $2ndd_f=2147483648$ 次乘加。此例说明二次项未必占据全部代价；以上只核算一层前向主要矩阵乘法，不含反向、Softmax、词表输出和实际硬件效率。
\end{example}

\begin{example}[三种缓存配置的比较]
取 $B=2$、$L=24$、缓存长度 $n=2048$、查询头数 $h=16$、单头宽度 $d_k=64$，每个数值两字节。MHA 的 $h_{\rm kv}=16$，式\eqref{eq:tr-kv-memory}给出 $402653184$ 字节，即 $384$ MiB；GQA 取 $h_{\rm kv}=4$ 时为 $96$ MiB；MQA 取 $h_{\rm kv}=1$ 时为 $24$ MiB，其中 $1\,\mathrm{MiB}=2^{20}$ 字节。这里比较的是理想缓存数据体积，不是模型全部显存，也不是速度比。若服务还有多个不同长度请求，应对各请求长度求和，再考虑分块分配与共享。
\end{example}

\begin{example}[逐词元平均与逐样本平均为什么不同]
某批次两个样本各有 $2$ 和 $8$ 个有效目标词元，平均负对数似然分别为 $1$ 和 $3$。逐样本等权平均得到 $(1+3)/2=2$，逐词元等权平均得到 $(2\times1+8\times3)/10=2.6$，对应困惑度分别约为 $7.389$ 和 $13.464$。两个数都来自可以明确定义的目标，但代表不同权重。若希望梯度累积等价于把所有有效词元放进一个大批次，就应累计损失总和并除以有效词元总数，不能简单平均不等长微批次的平均值。
\end{example}

\begin{example}[增量解码的掩码与位置核对]
假设缓存中已经有绝对位置 $0,1,2$，新输入对应位置 $3,4$。查询数为二，键数为五，允许矩阵应为
\[
\begin{pmatrix}1&1&1&1&0\\1&1&1&1&1\end{pmatrix}.
\]
第一行看不到位置四，第二行看得到全部前缀和自己；两个查询的位置编码应使用三和四，而不是重新从零开始。若使用 RoPE 且旧缓存键已经旋转，则只能旋转新键并追加，不能对拼接后的所有键再次旋转。这个例子可以同时检验缓存偏移、因果方向和旋转次数，是比大批量随机输入更容易解释的排错用例。
\end{example}

复习时可沿三个递进层次检查自己。第一层是闭卷画出数据流，写出 $Q,K,V$、分数、权重与输出的形状，说明掩码和标签移位。第二层是独立推出缩放方差、Softmax 梯度、RoPE 相对位置恒等式、参数数与缓存大小，并能说出各公式的假设。第三层是面对真实现象提出可区分原因的实验，例如缓存不一致时比较逻辑值和中间层，而不是只改变采样温度；训练损失异常低时检查信息泄漏，而不是直接庆祝收敛。

回答综合问题时，先明确任务与比较条件，再选择有解释力的量。例如“把上下文扩大四倍”在朴素稠密注意力矩阵上意味着约十六倍元素，在 KV 缓存上意味着四倍元素，在固定其余条件的单步注意力读取上约意味着四倍历史长度，但端到端延迟不必按其中任一倍数变化。把数学推导、实现约定和测量结果分开，是避免背诵结论却无法处理追问的关键。

\section{资料来源与使用边界}
\label{sec:tr-sources}

本文的手算例题、维度核算、梯度推导和问答按统一记号重新组织。文献用于追溯方法来源与适用条件，不把原论文某个硬件、数据集和训练配置下的提升幅度泛化为工程保证。框架接口可能随版本变化，使用参考代码迁移到实际项目时，应保存本地框架版本、模型配置和掩码约定。本文提供原理与面试准备材料，不声称已完成大模型训练、GPU 内核性能测试或招聘频率调查。

\begin{thebibliography}{99}
\bibitem{vaswani2017}
Vaswani A, Shazeer N, Parmar N, et al.
Attention Is All You Need. 2017.
\url{https://arxiv.org/abs/1706.03762}.

\bibitem{devlin2019}
Devlin J, Chang M W, Lee K, Toutanova K.
BERT: Pre-training of Deep Bidirectional Transformers for Language Understanding. 2018.
\url{https://arxiv.org/abs/1810.04805}.

\bibitem{raffel2020}
Raffel C, Shazeer N, Roberts A, et al.
Exploring the Limits of Transfer Learning with a Unified Text-to-Text Transformer. 2019.
\url{https://arxiv.org/abs/1910.10683}.

\bibitem{xiong2020}
Xiong R, Yang Y, He D, et al.
On Layer Normalization in the Transformer Architecture. ICML, 2020.
\url{https://proceedings.mlr.press/v119/xiong20b.html}.

\bibitem{zhang2019}
Zhang B, Sennrich R.
Root Mean Square Layer Normalization. 2019.
\url{https://arxiv.org/abs/1910.07467}.

\bibitem{shazeer2020}
Shazeer N. GLU Variants Improve Transformer. 2020.
\url{https://arxiv.org/abs/2002.05202}.

\bibitem{su2021}
Su J, Lu Y, Pan S, et al.
RoFormer: Enhanced Transformer with Rotary Position Embedding. 2021.
\url{https://arxiv.org/abs/2104.09864}.

\bibitem{press2021}
Press O, Smith N A, Lewis M.
Train Short, Test Long: Attention with Linear Biases Enables Input Length Extrapolation. 2021.
\url{https://arxiv.org/abs/2108.12409}.

\bibitem{ainslie2023}
Ainslie J, Lee-Thorp J, de Jong M, et al.
GQA: Training Generalized Multi-Query Transformer Models from Multi-Head Checkpoints. 2023.
\url{https://arxiv.org/abs/2305.13245}.

\bibitem{dao2022}
Dao T, Fu D Y, Ermon S, Rudra A, R\'e C.
FlashAttention: Fast and Memory-Efficient Exact Attention with IO-Awareness. 2022.
\url{https://arxiv.org/abs/2205.14135}.

\bibitem{kwon2023}
Kwon W, Li Z, Zhuang S, et al.
Efficient Memory Management for Large Language Model Serving with PagedAttention. 2023.
\url{https://arxiv.org/abs/2309.06180}.

\bibitem{hu2021}
Hu E J, Shen Y, Wallis P, et al.
LoRA: Low-Rank Adaptation of Large Language Models. 2021.
\url{https://arxiv.org/abs/2106.09685}.

\bibitem{fedus2021}
Fedus W, Zoph B, Shazeer N.
Switch Transformers: Scaling to Trillion Parameter Models with Simple and Efficient Sparsity. 2021.
\url{https://arxiv.org/abs/2101.03961}.

\bibitem{dosovitskiy2020}
Dosovitskiy A, Beyer L, Kolesnikov A, et al.
An Image is Worth 16x16 Words: Transformers for Image Recognition at Scale. 2020.
\url{https://arxiv.org/abs/2010.11929}.

\bibitem{pytorchsdpa}
PyTorch Contributors.
\texttt{scaled\_dot\_product\_attention}: official API documentation.
\url{https://docs.pytorch.org/docs/stable/generated/torch.nn.functional.scaled_dot_product_attention.html}.
访问日期：2026年10月3日；接口语义以使用版本为准。
\end{thebibliography}

\end{document}
